\insertImage{cfe76e97-852f-4871-9dda-e3a597c43a27.png}{Aquarelle d'une réunion d'équipe dans une entreprise libérée, avec des employés travaillant devant un tableau blanc et regardant ensemble un graphique d'entreprise sur le tableau blanc.}

\chapter{La performance et la rentabilité}

À un moment ou à un autre, la question tombe. En comité de direction, en conseil d'administration, au détour d'un rendez-vous bancaire : « D'accord, vos équipes sont épanouies. Mais est-ce que ça rapporte ? »

Et c'est une excellente question. Une entreprise libérée reste une entreprise : elle doit payer ses salaires, financer ses investissements, survivre à ses concurrents. La liberté ne dispense pas du compte de résultat. Ce chapitre regarde donc le modèle par son angle le moins romantique — et, croyez-moi, c'est un service à lui rendre : rien n'a fait plus de mal aux entreprises libérées que les promesses économiques qu'on a faites en leur nom.

\section{La mesure de la performance dans une entreprise libérée}

Premier problème, très concret : comment mesure-t-on la performance quand il n'y a plus de manager pour évaluer ?

Le réflexe classique — objectifs individuels fixés par le chef, entretien annuel, note — repose entièrement sur la hiérarchie. Supprimez-la, et deux tentations apparaissent. La première : ne plus rien mesurer du tout, au nom de la confiance. C'est une erreur, et une erreur qui se paie cher : sans repères partagés, personne ne sait si le collectif avance, et les problèmes se découvrent quand ils sont devenus des crises. La seconde tentation : garder les indicateurs d'avant en changeant juste les mots. « L'entretien annuel » devient « le moment feedback », et rien ne change. Les équipes reconnaissent l'ancien monde sous le déguisement en trente secondes.

Les entreprises libérées qui s'en sortent inventent une troisième voie : des mesures portées par les équipes elles-mêmes. Chez Valve, l'évaluation vient des pairs — ce sont les collègues qui apprécient la contribution de chacun, pas un supérieur\footnote{Valve Corporation (2012). Handbook for New Employees.}. Chez Buurtzorg, les équipes d'infirmières se pilotent sur des indicateurs qui ont du sens pour elles : la qualité des soins et la satisfaction des patients\footnote{Laloux, F. (2014). Reinventing Organizations: A Guide to Creating Organizations Inspired by the Next Stage of Human Consciousness. Nelson Parker.}. Le principe commun : l'indicateur appartient à ceux dont il mesure le travail. C'est toute la différence entre un tableau de bord et un outil de surveillance.

Et pour ceux qui douteraient qu'une organisation ainsi pilotée puisse être économiquement sérieuse, Buurtzorg fournit l'un des rares chiffres solides de toute la littérature : une étude d'Ernst \& Young a montré que ses équipes autonomes n'utilisaient qu'environ 40 \% des heures de soins autorisées par patient, contre environ 70 \% en moyenne chez leurs concurrents classiques — pour une qualité de soins supérieure et des patients plus vite autonomes\footnote{Gray, B. H., Sarnak, D. O., \& Burgers, J. S. (2015). Home Care by Self-Governing Nursing Teams: The Netherlands' Buurtzorg Model. The Commonwealth Fund.}. Moins d'encadrement, moins d'heures, de meilleurs soins. Le tableau de bord des infirmières a battu celui des contrôleurs de gestion.

Soyons honnêtes sur le coût : l'évaluation par les pairs est exigeante, parfois inconfortable, et elle demande une vraie maturité collective\footnote{Puranam, P., \& Håkonsson, D. D. (2015). Valve's way. Journal of Organization Design, 4(2), 2-4.}. Mais l'alternative — l'absence de mesure — est pire, parce qu'elle finit toujours de la même façon : par le retour brutal du contrôle, le jour où les chiffres déçoivent.

\section{Les risques de perte de compétitivité et de rentabilité}

Il faut prendre les sceptiques au sérieux, parce que leurs arguments ne sont pas absurdes.

Oui, la décision décentralisée peut être plus lente quand la coordination est mal outillée. Oui, l'investissement initial est réel : la formation, l'accompagnement, le temps passé à construire les nouveaux rituels — tout cela coûte avant de rapporter. Et oui, il existe un risque de confort : une entreprise focalisée sur son harmonie interne peut quitter des yeux son marché. Aucun de ces risques n'est une fatalité, mais les nier serait la meilleure façon de les subir.

Il y a aussi les résistances de l'écosystème, qu'on a déjà croisées : des clients, des investisseurs, des banquiers qui ne savent pas lire une entreprise sans organigramme et qui traduisent ce qu'ils ne comprennent pas en prime de risque.

Mais le tableau a un autre versant, documenté celui-là. FAVI a prospéré des décennies sur un marché — la sous-traitance automobile — réputé pour broyer ses fournisseurs, tout en restant en Picardie quand ses concurrents délocalisaient\footnote{Getz, I. (2009). Liberating Leadership: How the Initiative-Freeing Radical Organizational Form Has Been Successfully Adopted. California Management Review, 51(4), 32-58.}. Chrono Flex, le dépanneur nantais de flexibles hydrauliques, a engagé sa libération au sortir de la crise de 2009 et en a fait le socle de son rebond, raconté par son propre dirigeant\footnote{Gérard, A. (2017). Le patron qui ne voulait plus être chef. Flammarion.}. La liberté n'est pas l'ennemie de la performance. Elle en est une stratégie — exigeante, mais une stratégie.

Le vrai risque économique, à mes yeux, n'est ni la lenteur ni le coût. C'est la demi-mesure : une libération de façade qui paie tous les coûts de la transformation — le temps, la désorganisation, le scepticisme — sans jamais en toucher les bénéfices, parce que la confiance, elle, n'a jamais été réellement accordée.

\section{Les stratégies pour maintenir et améliorer la performance}

Alors, comment font celles qui durent ? Trois constantes reviennent, et vous allez voir qu'elles se ressemblent : ce sont les mêmes qui font réussir la transformation elle-même.

D'abord, la transparence économique. Dans les entreprises libérées qui performent, les chiffres ne sont pas réservés au comité de direction : les équipes connaissent les coûts, les marges, les enjeux. C'est la condition pour que l'autonomie produise de bonnes décisions — on ne peut pas décider intelligemment avec des informations qu'on n'a pas. La confiance sans les chiffres, c'est de la sympathie ; avec les chiffres, ça devient du pilotage.

Ensuite, l'investissement continu dans les compétences — on l'a vu au chapitre précédent, et il prend ici tout son sens économique : des équipes autonomes et compétentes décident vite et bien ; des équipes autonomes et perdues décident vite et mal.

Enfin, une vision qui sert de boussole. Chez FAVI, elle tenait en une obsession simple — l'amour du client, la qualité sans compromis — et c'est elle qui alignait les mini-usines sans qu'un directeur ait besoin de le faire\footnote{Getz, I. (2009). Liberating Leadership: How the Initiative-Freeing Radical Organizational Form Has Been Successfully Adopted. California Management Review, 51(4), 32-58.}. Une vision claire n'est pas un slogan sur un mur : c'est un critère de décision utilisable par n'importe qui, n'importe quel jour, sans demander la permission.

Transparence, compétence, boussole. Rien de magique — et c'est précisément le message de ce chapitre. La performance d'une entreprise libérée ne tombe pas du ciel avec la confiance. Elle se construit, elle se mesure, et elle s'entretient. Comme partout ailleurs. Juste pas par les mêmes personnes.
